\documentclass[11pt]{amsart}
\usepackage[mathscr]{eucal}
\usepackage[spanish, es-noshorthands]{babel}
\usepackage{times}
\usepackage{amsmath,amssymb,latexsym,amscd}
\usepackage{booktabs}   
\usepackage{hyperref}
\hypersetup{colorlinks=true,linkcolor=blue,citecolor=brown,linktocpage=true}
%\input xypic
%\xyoption{all}
\usepackage[all,cmtip]{xy}
\usepackage{fancyhdr}
\usepackage{mathalfa}
\usepackage{mathrsfs}
%\usepackage[sans]{dsfont}
\usepackage{upgreek}
%\usepackage{mathpazo}
\usepackage{tikz}
\usepackage{graphicx}
\usepackage{tikz-cd}
%%%%%%%%%%%%%%%%%%

\DeclareMathOperator{\op}{op}
\DeclareMathOperator{\spec}{spec}
\DeclareMathOperator{\pt}{pt}
\DeclareMathOperator{\tp}{tp}
\DeclareMathOperator{\fd}{fd}
\DeclareMathOperator{\id}{id}
\DeclareMathOperator{\Sp}{sp}
\DeclareMathOperator{\Ord}{Ord}
\DeclareMathOperator{\Frm}{Frm}
\DeclareMathOperator{\Top}{Top}
\newcommand{\niton}{\not\owns}

\theoremstyle{plain}
\newtheorem*{theorem*}{Teorema}
\newtheorem{thm}{Teorema}[section]
\newtheorem{cor}[thm]{Corolario}
\newtheorem{lem}[thm]{Lema}
\newtheorem{prop}[thm]{Proposición}

\theoremstyle{definition}
\newtheorem{dfn}[thm]{Definición}
\newtheorem{obs}[thm]{Observación}
\newtheorem{ej}[thm]{Ejemplo}

\renewcommand{\proofname}{Demostración}
  %%
    %\providecommand{\corollaryname}{Corollary}
  %\providecommand{\examplename}{Example}
  %\providecommand{\lemmaname}{Lemma}
  %\providecommand{\propositionname}{Proposition}
  %\providecommand{\remarkname}{Remark}
%\providecommand{\theoremname}{Theorem}
%\providecommand{\subexamplename}{Subexample}
  %\providecommand{\definitionname}{Definition}
\begin{document}
%\begin{frontmatter}

\title{Cosas probadas y preguntas abiertas}


\author{}
\email{}
\address{}


\address{}
%$\cortext[cor1]{Corresponding author}

%\begin{abstract}

%\end{abstract}

\maketitle
\section{Resultados probados}
\begin{prop}\label{fF}
Para $f=f^*\colon A\to B$ un morfismo de marcos y $G\in B^\wedge$, entonces $f_*G\in A^\wedge$.
\end{prop}

\begin{proof}
Por (\ref{Imagenfiltros}), $f_*G$ es un filtro en $A$. Necesitamos que $f_*G$ satisfaga la condición de filtro abierto. Sea $X\subseteq A$ tal que $\bigvee X\in f_*G$, con $X$ dirigido. Entonces
\[
Y=\{f(x)\mid x\in X\}
\] 
es dirigido y $f(\bigvee X)=\bigvee f[X]=\bigvee Y\in G$. Como $G$ es un filtro abierto y $f(\bigvee X)=\bigvee f[X]\in G$, existe $x\in X$ tal que $f(x)\in G$. Por lo tanto, $x\in f_*G$.
\end{proof}

En \cite[Lema 8.9 y Corolario 8.10]{simmons2004vietoris} muestran que el diagrama
\[\begin{tikzcd}
	A & A \\
	{\mathcal{O}S} & {\mathcal{O}S}
	\arrow["{f^\infty}", from=1-1, to=1-2]
	\arrow["{U_A}"', from=1-1, to=2-1]
	\arrow["{U_A}", from=1-2, to=2-2]
	\arrow["{F^\infty}"', from=2-1, to=2-2]
\end{tikzcd}\]
conmuta laxamente, es decir, $U_A\circ f^\infty\leq F^\infty \circ U_A$. En este diagrama $f^\infty$ y $F^\infty$ representan los
núcleos ajustados asociados a los filtros $F\in A^\wedge$ y $\nabla\in \mathcal{O}S^\wedge$, respectivamente.\\

Aquí probamos algo más general, ya que consideremos $j\in NA$, $F\in A_j^\wedge$ y $j_*F\in A^\wedge$ para producir el diagrama
\[\begin{tikzcd}
	A & A \\
	{A_j} & {A_j}
	\arrow["{f^\infty}", from=1-1, to=1-2]
	\arrow["j^*"', from=1-1, to=2-1]
	\arrow["j^*", from=1-2, to=2-2]
	\arrow["{\hat{f}^\infty}"', from=2-1, to=2-2]
\end{tikzcd}\]
donde $\hat{f}^\infty$ y $f^\infty$ son los núcleos ajustados asociados a los filtros $j_*F$ y $F$, respectivamente.\\

El siguiente lema es parte del análisis desarrollado en este trabajo y constituye una herramienta clave en la construcción de los diagramas conmutativos
porteriores.

\begin{lem}\label{f1f}
    Para $j=j^*$, $f$ y $\hat{f}$ como antes se cumple que $j\circ \hat{f}\leq f\circ j$.
\end{lem}
\begin{proof}
    Recordemos que, por (\ref{Imagenfiltros})
    \[
    \hat{f}=\dot{\bigvee}\{v_y\mid y\in j_*F\} \qquad f=\dot{\bigvee}\{v_{j(y)}\mid j(y)\in F\}. 
    \]
entonces para $a\in A$ se cumple
\[
v_y(a)=(y\succ a)\leq \hat{f}(a)\leq j(\hat{f}(a)).
\]

También, para todo $a, y\in A$, $(y\succ a)\wedge y=y\wedge a$ y
\[
\begin{split}
j((y\succ a)\wedge y)\leq j(a) & \Leftrightarrow j(y\succ a)\wedge j(y)\leq j(a)\\
& \Leftrightarrow j(y\succ a)\leq (j(y)\succ j(a)).
\end{split}
\]

Así 
\[
v_y(a)\leq j(\hat{f}(a))\leq (j(y)\succ j(a))=v_{j(y)}(j(a))\leq f(j(a)).
\]

Por lo tanto $j\circ \hat{f}\leq f\circ j$.
\end{proof}

Ahora probaremos lo anterior, pero para cualquier ordinal.

\begin{cor}\label{finftyf}
    Para $j$, $f$ y $\hat{f}$ como antes, se cumple que $j\circ \hat{f}^\alpha\leq f^\alpha\circ j$
\end{cor}

\begin{proof}
    Para $\alpha\in \Ord$ verificaremos que $j\circ \hat{f}^\alpha\leq f^\alpha\circ j$. Lo haremos con inducción transfinita.\\

    Si $\alpha=0$, es trivial.\\

    Para el paso de inducción, supongamos que para $\alpha$ se cumple. Luego
    \[
    j\circ \hat{f}^{\alpha+1}=j\circ \hat{f}\circ \hat{f}^{\alpha}\leq  f\circ j\circ \hat{f}^\alpha\leq f\circ f^\alpha\circ j=f^{\alpha+1}\circ j,
    \]
    donde la primera desigualdad es el Lema \ref{f1f} y la segunda es la hipótesis de inducción.\\

    Si $\lambda$ es un ordinal límite, entonces
	\[
	\hat{f}^\lambda=\bigvee\{\hat{f}^\alpha\mid \alpha<\lambda\}, \quad f^\lambda=\bigvee\{f^\alpha\mid \alpha<\lambda\} 
	\]
	y
	\[
		j\circ \hat{f}^\lambda=j\circ \bigvee_{\alpha<\lambda}\hat{f}^\alpha\leq\bigvee_{\alpha<\lambda}j\circ \hat{f}^\alpha.
	\]
	Por la hipótesis de inducción tenemos que 
	\[
	j\circ \hat{f}^\alpha\leq f^\alpha\circ j\Rightarrow \bigvee_{\alpha<\lambda}j\circ \hat{f}^\alpha\leq \bigvee_{\alpha<\lambda} f^\alpha\circ j.
	\]
	Por lo tanto $j\circ \hat{f}^\lambda\leq f^\lambda\circ j$.
\end{proof}

\begin{obs}\label{j monotono}
El Lema \ref{f1f} y el Corolario \ref{finftyf} también se cumplen si $j$ solo es monótono y preserva ínfimos finitos.
\end{obs}

Recordando que $v_F=f^\infty$ y $v_{j_*F}=\hat{f}^\infty$, el Corolario \ref{finftyf} nos dice que $j\circ v_{j_*F}\leq v_F\circ j$.\\

El objetivo ahora es comparar los cocientes compactos inducidos por $F$ y por $j_*F$, y construir un morfismo entre ellos que haga conmutar el diagrama correspondiente.\\

Si consideramos $j, v_F$ y $v_{j_*F}$ tenemos el siguiente diagrama
\begin{equation}\label{diagrama H}
    \begin{tikzcd}
	A && {A_{j_*F}} \\
	\\
	{A_j} && {A_F}
	\arrow["{(v_{j_*F})^*}", shift left, from=1-1, to=1-3]
	\arrow["j"', from=1-1, to=3-1]
	\arrow["H"', from=1-1, to=3-3]
	\arrow["{(v_{j_*F})_*}", shift left, from=1-3, to=1-1]
	\arrow["{(v_F)^*}"', shift right, from=3-1, to=3-3]
	\arrow["{(v_F)_*}"', shift right, from=3-3, to=3-1]
\end{tikzcd}
\end{equation}
donde $A_F$ y $A_{j_*F}$ son los cocientes compactos producidos por $v_F$ y $v_{j_*F}$, respectivamente.\\

Definimos el morfismo 
\[
H\colon A\to A_F, \qquad H=v_F\circ j.
\]

Queremos definir un morfismos de marcos entre los cocientes compactos $A_{j_*F}$ y $A_F$ de tal manera que 
el diagrama (\ref{diagrama H}) sea un cuadrado conmutativo.\\

Definimos
\[
h\colon A_{j_*F}\to A_j, \qquad h=H|_{A_{j_*F}},
\]
es decir, $h(x)=H(x)$ para todo $x\in A_{j_*F}$.\\

Notemos que, por el Corolario \ref{finftyf}, se cumple que $H(A_{j_*F})\subseteq A_F$, por lo que la restricción anterior 
está bien definida. Así, si $h=H_{\mid{A_{j_*F}}}$, el diagrama 
\[\begin{tikzcd}
	A & {A_{j_*F}} \\
	{A_j} & {A_F}
	\arrow["{v_{j_*F}}", from=1-1, to=1-2]
	\arrow["j"', from=1-1, to=2-1]
	\arrow["h", from=1-2, to=2-2]
	\arrow["{v_F}"', from=2-1, to=2-2]
\end{tikzcd}\]
conmuta.\\

Debemos verificar que $h$ es un morfismo de marcos. Por la definición de $h$, este es un $\wedge$-morfismo. Restaría comprobar que es un $\bigvee$-morfismo.\\

Los supremos en $A_{j_*F}$ y $A_F$ son calculados de manera diferente. De esta manera, denotamos por $\hat{\bigvee}$ al supremo en $A_{j_*F}$ y por $\tilde{\bigvee}$ al supremo en $A_F$. Por lo tanto
\[
\hat{\bigvee}=v_{j_*F}\circ \bigvee\quad\mbox{ y }\quad \tilde{\bigvee}=v_{F}\circ \bigvee,
\] 
es decir, para $X\subseteq A$, $Y\subseteq A_j$,
\[
	\hat{\bigvee}X=v_{j_*F}(\bigvee X)\quad\mbox{ y }\quad \tilde{\bigvee}Y=v_{F}(\bigvee Y).
\]

Como $H$ es un morfismo de marcos, $H\circ \bigvee=\tilde{\bigvee}\circ H$. Necesitamos algo similar para el morfismo $h$.

\begin{lem}\label{bigvee g}
$h\circ \hat{\bigvee}=\tilde{\bigvee}\circ h$.
\end{lem}

\begin{proof}
Basta probar que 
\[
h\circ \hat{\bigvee}\leq \tilde{\bigvee}\circ h
\]
ya que la otra desigualdad se obtiene por definición de los supremos en los cocientes.\\

Notemos que
\[
h\circ \hat{\bigvee}=H\circ v_{j_*F}\circ \bigvee=v_F\circ j\circ v_{j_*F}\circ \bigvee\leq v_F\circ v_F\circ j\circ \bigvee
\]
donde la desigualdad es el Corolario \ref{finftyf}. Además, por la idempotencia de $v_F$ 
\[
h\circ\hat{\bigvee}\leq v_F\circ j\circ \bigvee =H\circ \bigvee=\tilde{\bigvee}\circ H=\tilde{\bigvee}\circ h.
\]
Por lo tanto $h\circ\hat{\bigvee}=\tilde{\bigvee}\circ h$.
\end{proof}

La construcción anterior culmina en el siguiente resultado, que establece la conmutatividad del diagrama inducido por los cocientes compactos.
\begin{prop}\label{VFsquare}
El diagrama
\begin{equation}\label{Cuadrado}
    \begin{tikzcd}
	A & {A_{j_*F}} \\
	{A_j} & {A_F}
	\arrow["{v_{j_*F}}", from=1-1, to=1-2]
	\arrow["j"', from=1-1, to=2-1]
	\arrow["h", from=1-2, to=2-2]
	\arrow["{v_F}"', from=2-1, to=2-2]
\end{tikzcd}\end{equation}
es conmutativo.
\end{prop}

\begin{prop}\label{tidyquout}
    Si $A$ es un marco eficiente, entonces $A_j$ es un marco eficiente.
\end{prop}

\begin{proof}
Observemos primero que $F\subseteq j_*F$. Como $A$ es eficiente y $F\in A^\wedge$, se cumple que 
\[
x\in F\Rightarrow \hat{d}\vee x=1,
\]
donde $\hat{d}=d(\alpha)=f^\alpha(0)$.\\

Si $\hat{d}\leq d$, entonces $d\vee x=1$, para $d=d(\alpha)=f^\alpha(j(0))$.\\

Así, por el Corolario \ref{finftyf}
\[
\hat{d}=\hat{d}(\alpha)\leq j(\hat{d}(\alpha))=j(\hat{f}^\alpha(0))\leq f^\alpha(j(0))=d(\alpha)=d.
\]

Por lo tanto, si $x\in F$, entonces $d\vee x=1$ y $A_j$ es eficiente.
\end{proof}

Analizamos ahora la estabilidad de la eficiencia bajo coproductos binarios.
\begin{prop}\label{tidyplus}
Si $A_1, A_2\in \Frm$ son marcos eficientes, entonces $A_1\oplus A_2$ es un marco eficiente.
\end{prop}

\begin{proof}
Consideremos $A_1, A_2\in \Frm$ $\alpha$-eficientes. Por la construcción de $A_1\oplus A_2$ tenemos el siguiente diagrama
\[
\begin{tikzcd}
A_i \arrow[r, "\alpha_i"] \arrow[rrr, "\iota_i", bend left] & A_1\times A_2  \arrow[r, "\downarrow(\_)"] & \mathcal{L}(A_1\times A_2) \arrow[r, "\sim"] & A_1\oplus A_2 \arrow[lll, "(\iota_i)_*", bend left]
\end{tikzcd}
\]
donde $\alpha_1(a)=(a, 1)$ y $\alpha_2(b)=(1, b)$ para $a\in A_1$ y $b\in A_2$. Además, como $\iota_i\in \Frm$, $(\iota_i)_*$ es su adjunto derecho.\\

Consideremos $F\in (A_1\oplus A_2)^\wedge$ y, por la Proposición \ref{fF}, $(\iota_i)_*[F]\in A_i^\wedge$. Como $A_i$ es eficiente, para $x_i\in (\iota_i)_*F$ se cumple que $d_i\vee x_i=1$, donde 
\[
d_i=d_i(\alpha)=f_i^\alpha(0), \mbox{ con } f_i=\dot{\bigvee}\{v_y\mid y\in (\iota_i)_*F\}
\]
Por (\ref{Imagenfiltros}) se tiene que si $x_i\in (\iota_i)_*F$, entonces $\iota_i(x_i)\in F$, es decir.
\[
x_1\oplus 1, 1\oplus x_2\in F
\]
y al ser $F$ un filtro, se cumple que $x_1\oplus 1\wedge 1\oplus x_2=x_1\oplus x_2\in F$.\\

Además, para cada $i=1, 2$ tenemos el siguiente diagrama

\[\begin{tikzcd}
	{A_i} && {A_i} \\
	{A_1\oplus A_2} && {A_1\oplus A_2}
	\arrow["{f_i}", from=1-1, to=1-3]
	\arrow["{\iota_i}"', from=1-1, to=2-1]
	\arrow["{\iota_i}", from=1-3, to=2-3]
	\arrow["f_F"', from=2-1, to=2-3]
\end{tikzcd}\]
donde $f_F=\dot{\bigvee}\{v_{\iota_i(y)}\mid \iota_i(y)\in F\}$. Veamos que $\iota_i\circ f_i\leq f_F\circ \iota_i$.\\
    
Consideremos $a\in A_i$, entonces $v_y(a)\leq f_i(a)$ y $\iota_i(v_y(a))\leq \iota_i(f_i(a))$. Además,
\[
\begin{split}
\iota_i((y\succ a)\wedge y)=\iota_i(y\wedge a)\leq \iota_i(a)& \Leftrightarrow \iota_i(y\succ a)\wedge \iota_i(y)\leq \iota_i(a)\\
& \Leftrightarrow \iota_i(y\succ a)\leq (\iota_i(y)\succ \iota_i(a)).
\end{split}
\]
Luego 
\[
\iota_i(f_i(a))\leq (\iota_1(y)\succ \iota_i(a))=v_{\iota_i(y)}(\iota_i(a))\leq f_F(\iota_i(a)).
\]
Por lo tanto $\iota_i\circ f_i\leq f_F\circ \iota_i$. Verifiquemos que $\iota_i\circ f_i^\alpha\leq f_F^\alpha\circ \iota_i$ también es cierto. Por inducción transfinita 
\begin{itemize}
\item Si $\alpha=0$, es trivial.
\item Para el paso de inducción, supongamos que para $\alpha$ se cumple. Luego
\[
\iota_i\circ f_i^{\alpha+1}=\iota_i\circ f_i\circ f_i^{\alpha}\leq  f_F\circ \iota_i\circ f_i^\alpha\leq f_F\circ f^\alpha_F\circ \iota_i=f^{\alpha+1}_F\circ \iota_i,
\]
donde la primera desigualdad es lo demostrado antes y la segunda es la hipótesis de inducción.
\item Si $\lambda$ es un ordinal límite, entonces
\[
f_i^\lambda=\bigvee\{f_i^\alpha\mid \alpha<\lambda\}, \quad f^\lambda_F=\bigvee\{f^\alpha_F\mid \alpha<\lambda\}
\]
y
\[  
\iota_i\circ f_i^\lambda=\iota_i\circ \bigvee_{\alpha<\lambda}f_i^\alpha\leq\bigvee_{\alpha<\lambda}\iota_i\circ f_i^\alpha.
\]
Por la hipótesis de inducción tenemos que
\[
\iota_i\circ f_i^\alpha\leq f_F^\alpha\circ \iota_i\Rightarrow \bigvee_{\alpha<\lambda}\iota_i\circ f_i^\alpha\leq \bigvee_{\alpha<\lambda} f_F^\alpha\circ \iota_i.
\]
\end{itemize}
Así, evaluando en $0\in A_i$ se cumple que
\[
(\iota_i\circ f_i^\alpha)(0)\leq (f^\alpha_F\circ \iota_i)(0)\Rightarrow \iota_1(d_1)\leq f^\alpha_F(0\oplus 1)\;\mbox{ y }\;\iota_2(d_2)\leq f^\alpha_F(1\oplus 0).
\]
Luego 
\[
d_1\oplus d_2=d_1\oplus 1\wedge 1\oplus d_2\leq f^\alpha_F(0\oplus 1)\wedge f^\alpha_F(1\oplus 0)=f^\alpha_F(0\oplus 0). 
\]
De aquí que, para $d=d(\alpha)=f^\alpha_F(0\oplus 0)$, se cumple que
\[
(x_1\oplus x_2)\vee d_1\oplus d_2\leq (x_1\oplus x_2)\vee d,
\]
pero $x_1\oplus x_2\vee d_1\oplus d_2=(x_1\vee d_1)\oplus (x_2\vee d_2)=1\oplus 1$. Por lo tanto 
\[
x_1\oplus x_2\vee d=1\oplus 1,
\]
es decir, $A_1\oplus A_2$ es eficiente.
\end{proof}

\begin{cor}\label{Paso 3}
Si $A_1, A_2\in \Frm$ son marcos $\alpha$-eficiente y $\beta$-eficiente, respectivamente, entonces $A_1\oplus A_2$ es eficiente.
\end{cor}

\begin{proof}
Sin perdida de generalidad, consideremos $\gamma=\max\{\alpha, \beta\}$. Entonces $A_1$ y $A_2$ son $\gamma$-eficientes y por la Proposición \ref{tidyplus}, $A_1\oplus A_2$ es eficiente.
\end{proof}

Generalizamos el resultado anterior a coproductos arbitrarios.
\begin{prop}\label{coprodE}
La clase de marcos eficientes es cerrada bajo coproductos.
\end{prop}

\begin{proof}
Sean $\{A_i\}_{i\in I}$ una familia de marcos eficientes y $\{\alpha_i\}_{i\in I}\subseteq \Ord$ los grados de eficiencia de cada $A_i$. Sabemos que 
\[
\bigoplus_{i\in I} A_i\in \Frm \quad \mbox{ y }\quad (\iota_i\colon A_i\to \bigoplus_{i\in I} A_i)\in \Frm,
\]
entonces $(\iota_i)_*$ es el adjunto derecho de $\iota_i$. Además, para $F\in (\bigoplus_{i\in I} A_i)^\wedge$, por la Proposición \ref{fF}, $(\iota_i)_*F\in A_i^\wedge$ y por (\ref{Imagenfiltros}), 
\[
x_i\in (\iota_i)_*F\Leftrightarrow \iota_i(x_i)\in F.
\]
Sin perdida de generalidad, consideremos $\alpha=\sup\{\alpha_i\mid i\in I\}$, entonces cada $A_i$ es $\alpha$-eficiente. Por hipótesis, para $x_i\in (\iota_i)_*F$ se cumple que $x_i\vee d_i=1$, donde
\[
d_i=d_i(\alpha)=f_i^\alpha(0), \mbox{ con } f_i=\dot{\bigvee}\{v_y\mid y\in (\iota_i)_*F\}.
\]

\noindent
\textbf{Afirmación:} Para cada $x_i\in (\iota_i)_*F$ se cumple que $\langle x_i\rangle=\bigoplus \iota_i(x_i)\in F$.\\

Sabemos que $\bigoplus \iota_i(x_i)=\bigvee \iota_i(x_i)$. Luego, $\iota_i(x_i)\leq \bigvee \iota_j(x_j)$ para todo $i\in I$ y $j\in I$. 
Por lo tanto, al ser $F$ filtro, se cumple que $\bigvee \iota_i(x_I)\in F$.\\

Notemos que para cada $A_i$ tenemos el siguiente diagrama
\[\begin{tikzcd}
    {A_i} && {A_i} \\
    {\bigoplus_{i\in I} A_i} && {\bigoplus_{i\in I} A_i}
    \arrow["{f_i}", from=1-1, to=1-3]
    \arrow["{\iota_i}"', from=1-1, to=2-1]
    \arrow["{\iota_i}", from=1-3, to=2-3]
    \arrow["f_F"', from=2-1, to=2-3]
\end{tikzcd}\]
donde $f_F=\dot{\bigvee}\{v_{\iota_i(y)}\mid \iota_i(y)\in F\}$.\\

Si tomamos $a\in A_i$, de manera similar que en la prueba de la Proposición \ref{tidyplus}, se puede verificar que $\iota_i\circ f_i\leq f_F\circ \iota_i$ y por inducción transfinita, $\iota_i\circ f_i^\alpha\leq f_F^\alpha\circ \iota_i$. Evaluando en $0\in A_i$ se cumple que
\[
(\iota_i\circ f_i^\alpha)(0)\leq (f_F^\alpha\circ \iota_i)(0)\Leftrightarrow \iota_i(d_i)\leq f_F^\alpha(\langle 0\rangle=d(\alpha)=d.
\]
y $\langle d_i\rangle =\bigoplus \iota_i(d_i)\leq d$.\\

Luego, $\langle x_i\rangle \vee \langle d_i\rangle\leq \langle x_i\rangle \vee d$, pero 
\[
\langle x_i\rangle \vee \langle d_i\rangle=\langle x_i\vee d_i\rangle = \langle 1\rangle,
\]
es decir, $\langle 1\rangle \leq \langle x_i\rangle \vee d$ y por lo tanto $\langle x_i\rangle \vee d= \langle 1\rangle$. Así, por la afirmación anterior, $ \bigoplus_{i\in I} A_i$ es eficiente.
\end{proof} 

\begin{cor}
  La subcategoría de marcos eficientes es coreflexiva en $\Frm$.
\end{cor}

\begin{lem}\label{inducciontransfinita}
Para todo $\alpha\in \Ord$ y $X\in CS$, se cumple
\begin{equation}\label{eq:induccion}
\partial_F^\alpha(X)=\big(f_F^\alpha(X')\big)'.
\end{equation}
\end{lem}

\begin{proof}
Procedemos por inducción transfinita sobre $\alpha\in \Ord$.

\smallskip
\noindent\textbf{Caso base $\alpha=0$.}
Por definición, $\partial_F^0=\id$, es decir, $\partial_F^0(X)=X$ y $f_F^0=\id$, luego
\[
\partial_F^0(X)=X=(X'')=\big(f_F^0(X')\big)'.
\]

\smallskip
\noindent\textbf{Caso sucesor.}
Supongamos que la afirmación es cierta para $\alpha$ y probemos para $\alpha+1$.
Por definición de iteración,
\[
\partial_F^{\alpha+1}(X)=\partial_F\big(\partial_F^\alpha(X)\big).
\]
Además, por (\ref{eq:complementodual})
\[
\partial_F^{\alpha+1}(X)
=(f_F((\partial_F^\alpha(X))'))'.
\]
Por hipótesis inducción, $\partial_F^\alpha(X)=\big(f_F^\alpha(X')\big)'$ y sustituyendo obtenemos
\[
\partial_F^{\alpha+1}(X)=(f_F(f_F^\alpha(X')))'=\big(f_F^{\alpha+1}(X')\big)'
\]
que es lo que queríamos.

\smallskip
\noindent\textbf{Caso límite.}
Sea $\lambda$ un ordinal límite y supongamos que lo anterior se cumple para todo $\alpha<\lambda$.
Por definición de la iteración,
\[
\partial_F^\lambda(X)=\bigcap_{\alpha<\lambda}\partial_F^\alpha(X).
\]
Por De Morgan,
\[
\big(\partial_F^\lambda(X)\big)'
=\Big(\bigcap_{\alpha<\lambda}\partial_F^\alpha(X)\Big)'
=\bigcup_{\alpha<\lambda}\big(\partial_F^\alpha(X)\big)'.
\]
Usando la hipótesis de inducción, para todo $\alpha<\lambda$,
\[
\big(\partial_F^\alpha(X)\big)'=f_F^\alpha(X'),
\]
Luego,
\[
\big(\partial_F^\lambda(X)\big)'=\bigcup_{\alpha<\lambda} f_F^\alpha(X')=f_F^\lambda(X').
\]
Por lo tanto,
\[
\big(\partial_F^\lambda(X)\big)'=f_F^\lambda(X'),
\]
es decir,
\[
\partial_F^\lambda(X)=\big(f_F^\lambda(X')\big)'.
\]
y obtenemos lo que queríamos.
\end{proof}

Un caso particular del lema anterior es cuando $\alpha=\infty$.

\begin{cor}\label{cor:infinito}
Para $X\in \mathcal{O}S$ y $U\in \mathcal{O}S$ se cumple que
\begin{equation}\label{eq:infinito}
    \begin{split}
\partial_F^\infty(X)=(v_F(X'))' & \iff \partial_F^\infty(U')'=(v_F(U))\\
& \iff \partial_F^\infty(X)'=(v_F(X')).
    \end{split}
\end{equation}
\end{cor}

\begin{lem}\label{lem:auxiliar}
Si $A=\mathcal{O}S$ y $E\subseteq S$ se cumple que
\[
\mathcal{U}_*[E]\mathcal{U}=[E].
\]
\end{lem}

\begin{proof}
Como el marco $A$ es espacial se cumple que $\mathcal{U}$ es un isomorfismo. De aquí que si $U\in \mathcal{O}\pt A$ existe $V\in \pt A$ tal que $U=\mathcal{U}(V)$ y como $U\in \pt A$, se cumple que $U=V$. Luego
\[
\begin{split}
(\mathcal{U}_*[E]\mathcal{U})(U)&=\mathcal{U}_*[E](\mathcal{U}(U))=\mathcal{U}_*([E](U))\\
&=\mathcal{U}_*(E\cup U)^{\circ}=\bigcup\{W\in \mathcal{O}S\mid W\subseteq (E\cup U)^{\circ}\}\\
&=(E\cup U)^{\circ}=[E](U).
\end{split}
\]
Por lo tanto $\mathcal{U}_*[E]\mathcal{U}=[E]$.
\end{proof}

\begin{prop}\label{morfismo} .
    Para $F\in A^\wedge$, $Q\in\mathcal{Q}S$ y $\nabla=\nabla(Q)$, si $k\in [v_\nabla, w_\nabla]$, entonces $\mathcal{U}_*\circ k\circ\mathcal{U}\in [v_F, w_F]$.
\end{prop}

\begin{proof}
Si $k\in [v_\nabla, w_\nabla]$, entonces $\mathcal{U}_*\circ k\circ \mathcal{U}\in NA$. Además, para $F\in A^\wedge$ si $j\in [v_F, w_F]$, entonces $F=\nabla(j)$. De esta manera debemos probar que $F=\nabla(\mathcal{U}_*\circ k\circ \mathcal{U})$.\\

Por el Teorema de Hofmann-Mislove tenemos que
\[
x\in F\iff Q\subseteq \mathcal{U}(x)\iff \mathcal{U}(x)\in \nabla=\nabla(Q).
\]

Además, para todo $x\in A$
\[
(\mathcal{U}_*\circ k\circ\mathcal{U})(x)=\mathcal{U}_*(k(\mathcal{U}(x)))=\bigwedge(S\setminus k(\mathcal{U}(x))).
\]

Como $k\in [v_\nabla, w_\nabla]$, entonces $\nabla=\nabla(k)$, es decir, para $V\in \nabla$ se cumple que $k(V)=S$. Por lo tanto,
\[
\begin{split}
x\in F\iff \mathcal{U}(x)\in \nabla &\iff S\setminus k(\mathcal{U}(x))=\emptyset\\
&\iff \bigwedge(S\setminus k(\mathcal{U}(x)))=1\\
&\iff x\in \nabla(\mathcal{U}_*\circ k\circ \mathcal{U})
\end{split}
\]
\end{proof}

Lo anterior define un morfismo $\mho\colon [v_\nabla, w_\nabla]\to [v_F, w_F]$. Podría resultar interesante analizar que clase de morfismo es $\mho$. 
Por ejemplo, notemos que ocurre con los extremos del intervalo en $NA$.\\

Consideremos $k\in [v_\nabla, w_\nabla]$. Ya que $v_F$ es un núcleo ajustado, entonces se cumple que 
\begin{equation}\label{comparacionvF}
v_F\leq \mathcal{U}_*\circ k\circ \mathcal{U}. 
\end{equation}

De manera similar, dado que $w_F$ es el mayor elemento de su bloque, entonces
\begin{equation}\label{comparacionwF}
\mathcal{U}_*\circ k\circ \mathcal{U}\leq w_F.
\end{equation}

\begin{prop}\label{igualdadwF}
Para $F\in A^\wedge$ y $\nabla\in \mathcal{O}S^\wedge$ se cumple que 
\[
w_F=\mathcal{U}_*\circ w_\nabla\circ \mathcal{U}.
\]
\end{prop}

\begin{proof}
Por (\ref{comparacionwF}), debemos verificar que $w_F\leq \mathcal{U}_*\circ w_\nabla\circ \mathcal{U}$. Por la fórmula de la adjunción,
lo anterior es equivalente a 
\[
\mathcal{U}\circ w_F\leq w_\nabla\circ \mathcal{U},
\]
es decir, para todo $x\in A$ se cumple que 
\[
\mathcal{U}(w_F(x))\subseteq [M'](\mathcal{U}(x))=(M'\cup \mathcal{U}(x))^\circ.
\]
Ya que $\mathcal{U}(w_F(x))\in \mathcal{O}S$, es suficiente demostrar que $\mathcal{U}(w_F(x))\subseteq M'\cup \mathcal{U}(x)$.
A manera de contradicción, consideremos $y\in \mathcal{U}(w_F(x))$ tal que $y\notin M'$ y $y\notin \mathcal{U}(x)$. De aqui que
\[
y\in M\quad\text{ y }\quad x\leq y.
\]

Por hipótesis, $y\in \mathcal{U}(w_F(x))$, entonces 
\[
w_F(x)\nleq y\quad \text{ y }\quad w_F(x)=\bigwedge\{m\in M\mid x\leq m\},
\]
es decir, $w_F(x)\leq y$ lo cual es una contradicción.\\ 

Por lo tanto, $\mathcal{U}(w_F(x))\subseteq M'\cup \mathcal{U}(x)$.
\end{proof}


Un caso particular del diagrama (\ref{Cuadrado}) lo podemos obtener al considerar $j^*=\mathcal{U}_A$. Así,
\[\begin{tikzcd}
	A & {A_F} \\
	{\mathcal{O}S} & {\mathcal{O}S_\nabla}
	\arrow["{v_F}", from=1-1, to=1-2]
	\arrow["{\mathcal{U}_A}"', from=1-1, to=2-1]
	\arrow["h", from=1-2, to=2-2]
	\arrow["{v_\nabla}"', from=2-1, to=2-2]
\end{tikzcd}\]

Se puede probar que $\pt A_F=Q$, es decir, tenemos la reflexión espacial
\[
\mathcal{U}_{A_F}\colon A_F\to \mathcal{O}\pt A_F=\mathcal{O}Q
\]
y este es un cociente de $A$ a través de $v_F$.\\ 

Adicionalmente, tenemos el morfismo de marcos
\[\begin{tikzcd}
	A & {\mathcal{O}S} & {\mathcal{O}S_\nabla} & {\mathcal{O}Q}
	\arrow["{\mathcal{U}_A}", from=1-1, to=1-2]
	\arrow["{v_\nabla}", from=1-2, to=1-3]
	\arrow["{\mathcal{U}_{\mathcal{O}S_\nabla}}", from=1-3, to=1-4]
\end{tikzcd}\]
y es un cociente de $A$ a través de $v_\nabla$, pues se cumple que $\pt \mathcal{O}S_\nabla=Q$. Debido a que $Q\subseteq S$, ambos cocientes están dados por la asignación
\[
x\mapsto \mathcal{U}_A(x)\cap Q
\]
para $x\in A$. 

La construcción culmina en el siguiente diagrama, que sintetiza la interacción entre cocientes compactos, separación y levantamientos de núcleos.
\begin{equation}\label{Qcuadrado}
\begin{tikzcd}
	A && {A_F} & \\
	&&& {\mathcal{O}Q} \\
	{\mathcal{O}S} && {\mathcal{O}S_\nabla}
	\arrow["{v_F}", from=1-1, to=1-3]
	\arrow["{\mathcal{U}_A}"', from=1-1, to=3-1]
	\arrow["{\mathcal{U}_{A_F}}", from=1-3, to=2-4]
	\arrow["h"', from=1-3, to=3-3]
	\arrow["{v_\nabla}"', from=3-1, to=3-3]
	\arrow["{\mathcal{U}_{\mathcal{O}S_\nabla}}"', from=3-3, to=2-4]
\end{tikzcd}\end{equation}

Dado que el diagrama anterior esta codificado a través de $F\in A^\wedge$ y el $Q\in \mathcal{Q}S$ asociado, 
denominamos a (\ref{Qcuadrado}) como el \emph{$Q$-cuadrado}.

\begin{thm}\label{haus1}
Bajo las hipótesis del $Q$-cuadrado (\ref{Qcuadrado}). Si $A\in \Frm$ satisface la propiedad $\mathbf{(H)}$, entonces
\[
\mathcal{O}Q\cong \uparrow Q'.
\]
En otras palabras, el marco de abiertos del espacio de puntos de $A_F$ es isomorfo a un cociente compacto
y cerrado de un espacio Huasdorff.
\end{thm}

\begin{proof}
Por hipótesis, $A$ satisface $\mathbf{(H)}$. Como la propiedad $\mathbf{(H)}$ es cerrada bajo cocientes, $A_F$ y $\mathcal{O}S$ satisfacen $\mathbf{(H)}$.
Además, $\mathbf{(H)}$ es conservativa, es decir, $S=\pt A$ es un espacio $T_2$. Al ser $S$ $T_2$, si $Q\in \mathcal{Q}S$, entonces $Q\in \mathcal{C}S$ y, por la Proposición \ref{Propiedadesapilado}, $S$ 
es fuertemente apilado. Por la Definición \ref{def:apilado}, $v_\nabla=[Q']$.\\

Ya que $Q'\in \mathcal{O}S$, se cumple que $[Q']=u_{Q'}$ y este es complementado en $N\mathcal{O}S$. Además,
\[
\mathcal{O}S_{u_{Q'}}=\uparrow Q'=\mathcal{O}S_{v_\nabla}.
\]

Por \cite[Propocisión VI.3.3]{PicadoPultr} se cumple que $\mathcal{O}S_{v_\nabla}$ es espacial. De aquí que $g$ es un isomorfismo.\\

Por lo tanto, $\mathcal{O}\pt A_F\cong \mathcal{O}S_\nabla$, es decir, $\mathcal{O}Q\cong \uparrow Q'$.

\end{proof}

\begin{dfn}\label{def:apilado-marcos}
Sea $A\in \Frm$.
\begin{enumerate}
  \item Diremos que $A$ es fuertemente apilado si $\mathcal{U}_*[Q']\mathcal{U}=v_F$.
  \item Diremos que $A$ es apilado $(\mathcal{U}_*([Q'])\mathcal{U})(0)=v_F(0)$.
\end{enumerate}
\end{dfn}

Trivialmente tenemos que fuertemente apilado implica apilado. Nuestro objetivo es demostrar que las propiedades anteriores son conservativas, es decir, que un espacio $S$ es apilado (fuertemente apilado) si y sólo si su marco de abiertos $\mathcal{O}S$ es apilado (fuertemente apilado). 
Antes de hacer esto necesitamos el siguiente lema auxiliar.

\begin{prop}\label{apiladoconservativa}
Las nociones de fuertemente apilado y apilado son conservativas.
\end{prop}

\begin{proof}
Consideremos $S\in \Top$, $A=\mathcal{O}S$, $F\in A^\wedge$ y $Q\in \mathcal{Q}S$.\\

Primero, supongamos que $S$ es fuertemente apilado. De aquí que $v_F=[Q']$ y por el Lema \ref{lem:auxiliar} tenemos que $\mathcal{U}_*[Q']\mathcal{U}=[Q']=v_F$. Por lo tanto $\mathcal{O}S$ es fuertemente apilado.\\

Supongamos ahora que $S$ es apilado. Por lo anterior se cumple que 
\[
\overline{Q}=\partial_F^\infty(S)\iff \overline{Q}'=(\partial_F^{\infty}(S))'.
\]

Además, 
\[
\begin{split}
(\mathcal{U}_*[Q']\mathcal{U})(\emptyset)&=\mathcal{U}_*([Q'](\{V\in \pt\mathcal{O}S\mid \emptyset \not\subseteq V\}))\\
&=\mathcal{U}_*([Q'](\emptyset))=\mathcal{U}_*(Q')^{\circ}\\
&=\bigcup\{W\in \mathcal{O}S\mid W\subseteq (Q')^\circ\}\\
&=(Q')^{\circ}=\overline{Q}'.
\end{split}
\]
Por el Corolario \ref{cor:infinito} tenemos que $v_F(\emptyset)=(\partial_F^\infty(\emptyset'))'$. Por lo tanto 
\[
v_F(\emptyset)=(\mathcal{U}_*[Q']\mathcal{U})(\emptyset),
\]
es decir, $\mathcal{O}S$ es apilado.\\
\end{proof}

El análisis desarrollado en las subsecciones anteriores sugiere una posible conexión más profunda entre eficiencia y las nociones 
libres de puntos de apilado introducidas. En particular, surge la pregunta de si estas nociones permiten caracterizar la eficiencia
bajo hipótesis de separación adecuadas.\\

Las siguientes cuestiones forman parte del trabajo en curso y orientan las etapas posteriores de la investigación.\\

Recordemos que si un espacio es $T_1$ y fuertemente apilado, entonces se satisface que $w_\nabla=v_\nabla$. Adicionalmente, por la Proposición \ref{Propiedadesapilado}, se cumple que
si $S$ es apilado y $T_1$, entonces es sobrio.\\

Lo anterior nos hace cuestionarnos si las nociones en marcos tienen el mismo comportamiento, es decir, si $A$ es $T_1$ y fuertemente apilado, 
\[
\text{¿}v_F=w_F?
\]
Si lo anterior ocurre, ¿qué información nos proporciona con respecto al marco $A$? Además, se cumple que si $A$ es 
apilado y $T_1$, ¿$A$ es espacial?

\begin{prop}\label{v_F=w_F}
Sean $A\in \Frm$, $F\in A^\wedge$ y $Q\in \mathcal{Q}S$. Si $A$ es fuertemente apilado y $T_1$, entonces $v_F=w_F$. 
\end{prop}

\begin{proof}
Sabemos que para $F\in A^\wedge$, siempre se cumple que $v_F\leq w_F$.\\

Por fuertemente apilado y por la Proposición \ref{igualdadwF} tenemos que
\[
v_F=\mathcal{U}_*\circ [Q']\circ \mathcal{U}\quad \text{y}\quad w_F=\mathcal{U}_*\circ[M']\circ \mathcal{U},
\]
respectivamente. De aquí que
\[
\mathcal{U}_*\circ [Q']\circ \mathcal{U}\leq\mathcal{U}_*\circ[M']\circ \mathcal{U},
\]
Como $A$ es $T_1$, se cumple que $S$ es $T_1$, entonces $Q=M$. Por lo tanto
\[
\mathcal{U}_*\circ [Q']\circ \mathcal{U}=\mathcal{U}_*\circ[M']\circ \mathcal{U},
\]
es decir, $v_F=w_F$.
\end{proof}

Notemos que la Definición \ref{def:apilado-marcos} está dada para cualquier pareja $(F,Q)$
asociada por el Teorema de Hofmann--Mislove. En particular, si $q\in S=\pt A$, entonces
$\uparrow q\in \mathcal{Q}S$ y el filtro abierto asociado está dado por
\[
F=\{x\in A\mid \uparrow q\subseteq \mathcal{U}(x)\}.
\]

Si $A$ es apilado, se cumple que
\[
v_F(0)=(\mathcal{U}_*\circ[(\uparrow q)']\circ \mathcal{U})(0)
=\mathcal{U}_*(((\uparrow q)')^\circ).
\]

Por la fórmula del adjunto derecho,
\[
\mathcal{U}_*(((\uparrow q)')^\circ)
=\bigvee\{x\in A\mid \mathcal{U}(x)\subseteq (\uparrow q)'\}.
\]

Recordemos que
\[
\mathcal{U}(x)=\{p\in S\mid x\nleq p\},
\qquad
\mathcal{U}(x)'=\{p\in S\mid x\leq p\}.
\]

Por lo tanto,
\[
\mathcal{U}(x)\subseteq (\uparrow q)'
\quad\Longleftrightarrow\quad
\uparrow q\subseteq \mathcal{U}(x)'.
\]

Ahora bien,
\[
p\in \uparrow q
\quad\Longleftrightarrow\quad
q\sqsubseteq p
\quad\Longleftrightarrow\quad
p\leq q,
\]
pues $\sqsubseteq$ es el orden de especialización en $S$, el cual corresponde al orden
opuesto del marco $A$.\\

Así,
\[
\uparrow q\subseteq \mathcal{U}(x)'
\quad\Longleftrightarrow\quad
\forall p\in S,\ (p\leq q \Rightarrow x\leq p).
\]

De esta manera obtenemos
\[
v_F(0)
=\bigvee\{x\in A\mid \forall p\in S,\ (p\leq q \Rightarrow x\leq p)\}.
\]

Si además $A$ es $T_1$, entonces todo punto es máximo en $A$. En particular, si
$p\in S$ satisface $p\leq q$, necesariamente $p=q$. Por lo tanto la condición
anterior se reduce a $x\leq q$, y así

\[
v_F(0)
=\bigvee\{x\in A\mid x\leq q\}
=\bigvee \downarrow q
=q.
\]

Notemos que $v_F(0)\in A_F$ y $v_F(q)=v_F(0)=0_{A_F}$. Ya que $q$ es máximo en $A$, entonces $q$ es máximo en $A_F$.
Por lo tanto, $A_F=\{q<1\}$.\\

Ahora, $F=\nabla(v_F)$ y $x\in F$ si y solo si $v_F(x)=1$. De esta manera, bajo las hipótesis anteriores, se cumple que 
\[
v_F(x)= \left\{ \begin{array}{lcc} 1, & \text{si}  &  x\nleq q\\ \\ q, & \text{si} & x\leq q 
\end{array} \right.=  \left\{ \begin{array}{lcc} 1, & \text{si}  &  x\in F\\ \\ q, & \text{si} & x\not\in F \end{array} \right.,
\]
es decir, la pareja $(F, \uparrow q)$ proporciona un filtro abierto que permite distinguir 
cuando un elemento de $A$ está en el filtro o no.

\begin{prop}\label{EficienteFapilado}
  Si $A\in \Frm$ es eficiente y el diagrama 

\[\begin{tikzcd}
	A & A \\
	{\mathcal{O}S} & {\mathcal{O}S}
	\arrow["{v_F}", from=1-1, to=1-2]
	\arrow["{\mathcal{U}}"', from=1-1, to=2-1]
	\arrow["{\mathcal{U}}", from=1-2, to=2-2]
	\arrow["{v_\nabla}"', from=2-1, to=2-2]
\end{tikzcd}\]
conmuta, entonces $A$ es fuertemente apilado.
\end{prop}

\begin{proof}
Sean $A\in\Frm$, $S=\pt A$ y $(F, Q)$ un par HM. Por hipótesis, $A$ es eficiente lo cual implica lo siguiente:
\[
A\text{ eficiente }\Rightarrow \mathcal{O}S\text{ eficiente }\Rightarrow S \text{ apilado y empaquetado.}
\]
De aquí que si $Q\in \mathcal{Q}S$, entonces $Q\in \mathcal{C}S$ y $Q'\in \mathcal{O}S$. Luego, al ser $S$ apilado tenemos que
$\overline{Q}=Q(\infty)$, es decir, $Q'=v_\nabla(\emptyset)$ y al ser $\mathcal{O}S$ eficiente, $[Q']=v_\nabla$.\\
  

Además, si $v_\nabla\circ \mathcal{U}=\mathcal{U}\circ v_F$, entonces
\[
v_\nabla\circ \mathcal{U}\leq \mathcal{U}\circ v_F\quad \text{ y } \mathcal{U}\circ v_F\leq v_\nabla\circ \mathcal{U},
\]
es decirc 
\[
[Q']\circ \mathcal{U}\leq \mathcal{U}\circ v_F\quad \text{ y } \mathcal{U}\circ v_F\leq [Q']\circ \mathcal{U}.
\]
Por propiedades del adjunto derecho 
\[
\mathcal{U}_*\circ [Q']\circ \mathcal{U}\leq v_F\quad \text{ y } v_F\leq\mathcal{U}_*\circ [Q']\circ \mathcal{U}.
\]
Por lo tanto $v_F=\mathcal{U}_*\circ [Q']\circ \mathcal{U}$, es decir, $A$ es fuertemente apilado.
\end{proof}

Un corolario que surge de la prueba anterior es el siguiente

\begin{cor}
Si $\mathcal{O}S$ es eficiente y $S$ es empaquetado, entonces $S$ es fuertemente apilado.
\end{cor}

\subsection{El marco $[0,1]$}
Siguiendo la idea presentada en \cite{simmons2006regularity} (o, de manera alternativa, en \cite{simmons2004vietoris}), consideremos el marco linealmente ordenado 
\[
A=[0,1].
\] 
En este caso se cumplen las siguientes observaciones.
\begin{enumerate}
\item Para todo $a,b\in A$ la implicación $(a\succ b)$ está dada por
\[
(a\succ b)=\begin{cases}
1, & \text{si } a\leq b,\\
b, & \text{si } a>b.
\end{cases}
\]
\item Para $a, x\in A$, los núcleos distinguidos $u_a, v_a, w_a$ están dados por
\[
u_a(x)=\begin{cases}
  x & \text{si }a\leq x,\\
  a & \text{si }x< a,
\end{cases}, \quad v_a(x)=\begin{cases}
1 & \text{si } x\geq a,\\
a & \text{si } x<a,
\end{cases} \quad w_a(x)=\begin{cases}
1 & \text{si } x>a,\\
a & \text{si } x\leq a.
\end{cases}
\]

\item Si $F\subseteq A$ es un filtro, entonces $F$ es siempre admisible y tiene la forma
\[
F=\nabla(w_a)=(a, 1]\quad \mbox{ o }\quad F=\nabla(v_a)=[a, 1]
\]
para algún $a\in A$. En particular, $F\in A^\wedge$ si y solo si $F=(a, 1]$.

\item Si $F\in A^\wedge$, el núcleo ajustado asociado $v_F$ se calcula como
\begin{equation}\label{eq:lm}
v_F(x)=l_m(x)=\begin{cases}
1 & \text{si } x> m,\\
x & \text{si } x\leq m.
\end{cases}
\end{equation}
donde $l_m=v_m\wedge w_m$ y $m=\bigwedge F$.
\item Para $S=\pt A$ se cumple que $S=[0,1)$ y los abiertos en $\mathcal{O}S$ son de la forma
\[
\mathcal{U}_A(x)=\{p\in S\mid p<x\}=[0,x).
\]
\item El marco $A$ no es $T_1$, pues ningún $p\in S$ es máximo. En consecuencia, $A$ no satisface ninguna propiedad de separación más fuerte en $\Frm$.
\item Para $F\in A^\wedge$ el conjunto compacto saturado asociado $Q\in \mathcal{Q}S$ es de la forma $Q=[0,m]$.
\item Se cumple que $[Q']=v_F$; por lo tanto, $S$ es fuertemente apilado.
\item Ya que $[0, x)\in \mathcal{O}S$, entonces $[x, 1]\in \mathcal{C}S$ y $Q=[0, m]$. Por lo tanto, los subconjuntos compactos no son cerrados; es decir, $S$ no es empaquetado.
\item El marco $A$ no es eficiente.

\end{enumerate}

Todo lo anterior resume las principales características del marco $[0,1]$. En lo que sigue, realizaremos el análisis de las construcciones de parches tanto de 
$A$ como de su espacio de puntos $S$.\\

Recordemos que para $A\in \Frm$ y $S\in \Top$ tenemos 
\[
\text{Pbase}=\{u_a\wedge v_F\mid a\in A, F\in A^\wedge\}\quad \text{y}\quad \text{pbase}=\{U\cap Q'\mid U\in \mathcal{O}S, Q\in \mathcal{Q}S\}
\]
donde Pbase genera al marco de parches de $A$ y pbase genera la topología del espacio de parches de $S$. Por \eqref{eq:lm} tenemos que 
\[
\text{Pbase}=\{u_a\wedge l_m\mid a\in [0,1], m\in [0,1)\}.
\]

Para $x\in A$  se cumple
\[
(u_a\wedge l_m)(x)=u_a(x)\wedge l_m(x)..
\]

Analizamos por casos la relación entre $x$ y $m$.
\begin{description}
\item[Caso 1] Si $x<m$, entonces $(u_a\wedge l_m)(x)=u_a(x)\wedge x=x$.
\item[Caso 2] Si $x=m$, entonces $(u_a\wedge l_m)(x)=u_a(m)\wedge m=x$.
\item[Caso 3] Si $x>m$, entonces $(u_a\wedge l_m)(x)=u_a(x)\wedge 1=u_a(x)$.
\end{description}

En los primeros 2 casos el valor de $a$ no afecta la evaluación. La dependencia esencial aparece cuando Además $u_a\wedge l_m=\id$, es decir, los generadores básicos de la 
$m<x$, donde interviene la relación entre $a$ y $m$.
\begin{description}
\item[Caso 3.1] Si $m<x<a$, entonces $(u_a\wedge l_m)(x)=a$.
\item[Caso 3.2] Si $m<a\leq x$, entonces $(u_a\wedge l_m)(x)=x$.
\end{description}

De este análisis se deduce que, $u_a\wedge l_m=\id$ si y solo si $m<a$. Por lo tanto, lo generadores no triviales de Pbase corresponden exactamente a los pares
$(m, a)$ con $m<a$.\\

En consecuencia, el marco $PA$ queda codificado por los intervalos abiertos de la forma $(m,a)$, donde $m\in \pt A$ y $a\in A$.

Queremos identificar si en nuestro ejemplo, $PA$ satisface alguna propiedad de separación. Para hacerlo, ocuparemos recordar información respecto a los parches.\\

Si $A\in\Frm$ es un marco arbitrario y $S$ su espacio de puntos, podemos construir el siguiente diagrama conmutativo que relaciona ambas construcciones de parches.
\[\begin{tikzcd}
	A & PA & NA \\
	{\mathcal{O}S} & {P\mathcal{O}S} & {N\mathcal{O}S} \\
	& {\mathcal{O}^pS} & {\mathcal{O}^fS}
	\arrow[from=1-1, to=1-2]
	\arrow["{\mathcal{U}_A}"', from=1-1, to=2-1]
	\arrow[hook, from=1-2, to=1-3]
	\arrow["{P(\mathcal{U}_A)}", from=1-2, to=2-2]
	\arrow["{N(\mathcal{U}_A)}", from=1-3, to=2-3]
	\arrow[from=2-1, to=2-2]
	\arrow[hook, from=2-1, to=3-2]
	\arrow[hook, from=2-2, to=2-3]
	\arrow["\pi", from=2-2, to=3-2]
	\arrow["\sigma", from=2-3, to=3-3]
	\arrow[hook, from=3-2, to=3-3]
\end{tikzcd}\]
donde $\mathcal{O}^fS$ es la topología de Skula (o topología frontal) del espacio $S$. En \cite{sexton2006point} lo llaman el 
\emph{diagrama del ensamble de parches} y más detalles sobre este también pueden ser encontrados
en \cite{sexton2003point}.\\

En nuestro caso, $\mathcal{U}_A$ es un isomorfismo. Además, bajo $\mathcal{U}_A$, la construcción de parches resulta ser funtorial y, por lo tanto,
$PA\cong P\mathcal{O}S$. De manera similar, $NA\cong N\mathcal{O}S$.
Para describir explícitamente $P\mathcal{O}S$ y $N\mathcal{O}S$ utilizamos el siguiente resultado, que puede consultarse en \cite{AvilaBezhaniMorandiZaldivar2020}.

\begin{prop}\label{Ensamble espacial}
Sea $S$ un conjunto parcialmente ordenado.
\begin{enumerate}
  \item Si $S$ no tiene anticadenas infinitas, entonces $N\mathcal{O}S$ es espacial.
  \item Si $S$ es totalmente ordenado, entonces $N\mathcal{O}S$ es espacial.
\end{enumerate}
\end{prop}

Así, por la Proposición \ref{Ensamble espacial} tenemos que $N\mathcal{O}S$ es espacial y por lo tanto $N\mathcal{O}S\cong \mathcal{O}^fS$. Juntando toda esta información 
el diagrama del ensamble de parches en nuestro caso particular queda como sigue.

\[\begin{tikzcd}
	A & PA & NA \\
	{\mathcal{O}S} & {\mathcal{O}^pS} & {\mathcal{O}^fS}
	\arrow[from=1-1, to=1-2]
	\arrow["\cong"', from=1-1, to=2-1]
	\arrow[from=1-2, to=1-3]
	\arrow["\cong", from=1-2, to=2-2]
	\arrow["\cong", from=1-3, to=2-3]
	\arrow[from=2-1, to=2-2]
	\arrow[from=2-2, to=2-3]
\end{tikzcd}\]
Por lo tanto, si queremos estudiar al marco $PA$ podemos estudiar a la topología de parches $\mathcal{O}^pS$. Recordemos que 
\[
\text{pbase}=\{U\cap Q'\mid U\in \mathcal{O}S, Q\in \mathcal{Q}S\}.
\]
En nuestro caso,
\[
\text{pbase}=\{[0, a)\cap [m, 1)\mid a\in [0,1], m\in [0,1)\}=\{(m,a)\}.
\]

Como la topología generada por los intervalos $(m,a)$ coincide con la topología usual en $[0,1)$, el espacio $^pS$ es Hausdorff. En consecuencia, $\mathcal{O}^pS\cong PA$ satisface $\mathbf{(H)}$, aun cuando $A$ no es $T_1$ ni eficiente.

\begin{dfn}\label{KC y KCH}
Sea $A\in \Frm$ decimos que $A$ es:
\begin{enumerate}
    \item $KC$ si cada cociente compacto es cerrado.

    \item \emph{Hausdorff cerrado compacto} (o $KCH$ de manera abreviada) si cada cociente compacto de $A$ es cerrado y Hausdorff.
\end{enumerate}
\end{dfn}

La Definición \ref{KC y KCH} nos dice que si $j\in NA$, entonces $\nabla(j)\in A^\wedge$ y $j=u_a$ para algún $a\in A$. De manera adicional, para $2)$ pedimos que $A_j$ cumpla la propiedad $\mathbf{(H)}$.\\

\begin{prop}\label{KCquout}
    Si $A\in \Frm$ comple $KC$, entonces $A_j$ cumple $KC$ para cada $j\in N(A).$
\end{prop}

\begin{proof} %HAY que redactyar mejor la prueba%
Consideremos $k\in NA_j$ tal que $(A_j)_k$ es compacto. 
Como cualquier filtro abierto es admisible, tenemos que $\nabla(k)\in A_j^\wedge$ 
y por la Proposición \ref{fF} $j_*\nabla(K)\in A^\wedge$.\\

Sea $l=j_*\circ  k\circ j^*\in NA$, entonces $A_l$ es un cociente compacto de $A$ y existe $a\in A$ tal que $l=u_a$. Así
\[
\begin{tikzcd}
	A \arrow[r, "j^*"'] \arrow[rrr, "l", bend left] & A_j \arrow[r, "k"'] & (A_j)_k \arrow[r, "j_*"'] & A_j\subseteq A
	\end{tikzcd}\]
y $a\vee x=k(j(x))$. Por lo tanto, si $x=a$, $k(j(x))=a$.\\

Necesitamos que $k=u_b$ para algún $b\in A_j$. Para $x\in A_j$ y $b=j(a)$
\[
\begin{split}
u_b(x)= b\vee x= b\vee j(x)& =j(j(a)\vee j(x))\\
& =j(k(j(a))\vee x)\\
& =j(u_a(x))\\
& =j(k(x))\\	
&=k(x).
\end{split}
\]
Por lo tanto $u_b=k$.
\end{proof}

\begin{prop}\label{KCT1}
Si $A$ es un marco $KC$, entonces $A$ es un marco $T_1$.
\end{prop}

\begin{proof}
Sean $p\in \pt A$ y $a\in A$ tales que $p\leq a\leq 1$. Consideremos 
\[
w_p(x)=\left\{\begin{array}{lcc}
1 & \mbox{ si } & x\nleq p\\
\\
p & \mbox{ si } & x\leq p
\end{array}\right.
\]
para $x\in A$. $P=\nabla(w_p)=\{x\in A\mid x\nleq p\}$ es un filtro completamente primo (en particular, $P\in A^\wedge$). Como $A$ es $KC$, entonces $A_{w_p}$ es un cociente compacto cerrado. Así $u_p=w_p$ y
\[
u_p(a)=a\quad \mbox{and}\quad w_p(a)=1.
\]
es decir, $a=1$. Por lo tanto $p$ es máximo. 
\end{proof}

Nuestro objetivo es estudiar los núcleos que producen cociente compacto. De manera particular, estamos interesados en 
aquellos núcleos que producen cociente compacto y cerrado.\\

Lo presentado en esta sección es una generalización de lo hecho por Escardó en \cite{escardo2006compactly}. En este caso, el lugar de utilizar la propiedad
$\mathbf{(fH)}$, usamos la eficiencia (o $KC$ cuando sea necesario).

\begin{dfn}\label{Definicion2.1}
Para $A\in \Frm$ y $j\in NA$, decimos que $j$ es $kq$ (por cociente compacto) si $A_j$ es compacto.  
\end{dfn}

Denotamos por 
\[
\mathfrak{K}A=\{j\in NA\mid j \mbox{ es } kq\}.
\]

Sabemos que los $u$-núcleos producen cocientes cerrados. De esta manera, un $u$-núcleo produce cociente compacto si $u_{\bullet}\in \mathfrak{K}A$ para $\bullet\in A$. Consideremos los subconjuntos
\[
\mathfrak{C}A=\{u_\bullet\in NA\mid u_\bullet\in \mathfrak{K}A\}\quad \mbox{ y }\quad \mathfrak{c}A=\{c\in A\mid u_c\in \mathfrak{K}A\}.
\]
donde $\mathfrak{C}A\cong \mathfrak{c}A$. El uso de cada uno de ellos depende del enfoque que ocupemos, ya sea como elementos del marco principal o como núcleos.

Para $j, k\in NA$ tenemos la relación de equivalencia dada por 
\[
j\sim k \Leftrightarrow \nabla(j)=\nabla(k),
\]
donde $\nabla( \_ )$ es el filtro de admisibilidad correspondiente al respectivo núcleos.\\

La clase de cada núcleo define un bloque en $NA$ y se puede demostrar que cada uno de estos bloques tiene menor elemento. Al menor elemento del bloque se le conoce como núcleo ajustado y todo núcleo ajustado tiene la forma
\[
f=\dot\bigvee \{v_a\mid a\in F\}
\]
donde $F$ es un filtro en $A$ y $\dot{\bigvee}$ es el supremo puntual.

\begin{lem}\label{cn ajustado}
    Consideremos $j\in NA$ y supongamos que $j$ es ajustado. Entonces
    \[
    j\leq k\Leftrightarrow \nabla(j)\subseteq \nabla(k)
    \]
    para todo $k\in NA$.
\end{lem}

Si $F\in A^\wedge$, $F=\nabla(j)$ para algún $j$ en $NA$, el bloque de $j$ siempre tiene un mayor y un menor elemento. Denotamos al menor elemento del bloque por  $v_F=f^\infty$.\\

Por último, si el marco $A$ es eficiente, se cumple que $v_F=u_d$, donde $d=v_F(0)$. Notemos que la eficiencia proporciona un cociente compacto y cerrado (el respectivo $v_F$), 
para cada $F\in A^\wedge$.\\

Con todo lo anterior tenemos la siguiente relación entre los distintos núcleos mencionado hasta este momento
\[
\mathfrak{f}\mathfrak{C}A\subseteq \mathfrak{C}A\subseteq \mathfrak{K}A\subseteq NA
\]
donde $\mathfrak{f}\mathfrak{C}A$ es el conjunto de núcleos $v_F=u_d$ que se obtienen cuando un marco es eficiente. Además, con respecto a los marcos se cumple que
\[
KCH\Rightarrow KC\Rightarrow \mbox{Eficiente}.
\]

\begin{lem}\label{Lema2.2}
    Para $c\in A$ las siguientes son equivalentes:
    \begin{enumerate}
        \item $c\in \mathfrak{c}A$.
        \item $A_{u_c}$ es un marco compacto.
        \item $\nabla(u_c)\in A^\wedge$.
    \end{enumerate}
\end{lem}

\begin{proof}
    $1)\Rightarrow 2)$ se cumple por caracterización de marcos compactos. Si se cumple $3)$, entonces $c\in \mathfrak{c}A$. Por último, $2) \Leftrightarrow 3)$ es cierto por la caracterización de cocientes compactos y filtros admisibles. 
\end{proof}

\begin{lem}\label{Lema2.3}
    Lo siguiente se cumple:
    \begin{enumerate}
        \item $\mathfrak{c}A$ es una sección superior.
        \item Si $X\subseteq \mathfrak{c}A$, entonces $\bigvee X\in \mathfrak{c}A$ y si $c, c'\in \mathfrak{c}A$ entonces $(c\succ c')\in \mathfrak{c}A$.
        \item Si $c, c'\in \mathfrak{c}A$, entonces $c\wedge c'\in \mathfrak{c}A$.
        \item $\mathfrak{c}A\subseteq A$ es un premarco con implicación.
        \item $\mathfrak{c}A$ es un premarco compacto.
    \end{enumerate}
\end{lem}

\begin{proof}
    \begin{enumerate}
        \item Sean $c\leq c'$ tal que $c\in \mathfrak{c}A$, entonces $\nabla(u_c)\subseteq \nabla(u_{c'})$. Consideremos $X\subseteq A$ dirigido con $\bigvee X\in \nabla(u_{c'})$. Debemos probar que 
        \[
        \nabla(u_{c'})\cap X\neq \emptyset.
        \]
        
        Notemos que 
        \[
        c'\vee (\bigvee X)=1\Rightarrow c\vee (c'\vee (\bigvee X))=1\Rightarrow c'\vee (\bigvee X)\in \nabla(u_c). 
        \]
        y el conjunto $Y=\{c'\vee x\mid x\in X\}$ es dirigido tal que $\bigvee Y\in \nabla(u_c)$. Al ser $\nabla(u_c)$ abierto, se tiene que existe $y\in Y$ tal que $y\in \nabla(u_c)$ con $y=c'\vee x$. Además, $y\in \nabla(u_{c'})$, es decir,
        \[
        y\vee c'=(c'\vee x)\vee c'= c'\vee x=1
        \]
        Por lo tanto $x\in \nabla(u_{c'})\in A^\wedge$ y así $c'\in \mathfrak{c}A$.
        \item Se cumple por ser sección superior.
        \item Si $c, c'\in \mathfrak{c}A$, entonces $u_c, u_{c'}\in \mathfrak{K}A$. De esta manera, debemos verificar que $u_{c\wedge c'}\in \mathfrak{K}A$. Por propiedades de los $u$-núcleos $u_{c\wedge c'}=u_c\wedge u_{c'}$, entonces
        \[
            \nabla(u_{c\wedge c'})=\nabla(u_c\wedge u_{c'})=\nabla(u_c)\cap \nabla(u_{c'}).
        \]
        Consideremos $X\subseteq A$ dirigido tal que $\bigvee X\in \nabla(u_{c\wedge c'})$. De aquí que $\bigvee X\in \nabla(u_c)\cap \nabla(u_{c'})$. Como ambos son filtros abiertos, existe $x\in X$ tal que $x\in \nabla(u_c)$ y $x\in \nabla(u_{c'})$, es decir, $x\in \nabla(u_c)\cap \nabla(u_{c'})$. Por lo tanto $\nabla(u_{c\wedge c'})\in A^\wedge$, es decir, $c\wedge c'\in \mathfrak{C}A$.
        \item Es consecuencia de $1)$, $2)$ y $3)$.
        \item Sea $X$ un conjunto dirigido de elementos en $\mathfrak{c}A$ tal que $\bigvee X=1$ y consideremos $c\in \mathfrak{c}A$. Como $u_c$ es $kq$, entonces $A_{u_c}$ es compacto y así $1$ es compacto en $A_{u_c}$, es decir, para
        \[
        \bigvee X=\bigvee^{u_c}X=1
        \]
        existe $x\in X$ tal que $x=1$ y $1\in A_{u_c}$.  Por lo tanto, $\mathfrak{c}A$ es compacto.
    \end{enumerate}
\end{proof}

\begin{lem}\label{Lema2.8y2.9}
    Sea $A\in \Frm$ y $\alpha_A\colon \mathfrak{c}A\to A^\wedge$ la función definida por
    \[
    \alpha_A(c)=\{x\in A\mid c\vee x=1\}=\nabla(u_c).
    \]
    Lo siguiente se cumple:
    \begin{enumerate}
        \item $\alpha_A$ es un morfismo de premarcos.
        \item Si $A$ es eficiente, $\alpha_A$ es sobreyectiva.
    \end{enumerate}
\end{lem}

\begin{proof}
    \begin{enumerate}
        \item Veamos primero que $\alpha_A$ preserva la estructura de un premarco.
    \begin{enumerate}
        \item Si $c\leq c'$,  entonces $\alpha_A(c)=\nabla(u_c)\subseteq \alpha_A(c')=\nabla(u_{c'})$.
        \item Consideremos $1\in \mathfrak{c}A$, entonces 
        \[
        \alpha_A(1)=\nabla(\tp)=A,
        \]
        de aquí que, $\alpha_A$ preserva el mayor elemento.
        \item Sean $c_1, c_2\in \mathfrak{c}A$ y $X\subseteq \mathfrak{c}A$ dirigido. De aquí que
        \[
        \begin{split}
            \alpha_A(c_1\wedge c_2)&=\{x\in A\mid (c_1\wedge c_2)\vee x=1\}\\
            &=\{x\in A\mid (c_1\vee x)\wedge (c_2\vee x)=1\}\\
            &=\{x\in A\mid c_1\vee x=1 \mbox{ y }c_1\vee x=1\}\\
            &=\nabla(u_{c_1})\cap \nabla(u_{c_2})\\
            &=\alpha_A(c_1)\cap \alpha_A(c_2).
        \end{split}
        \]
        y
        \[
        \begin{split}
        \alpha_A(\bigvee X)&=\{x\in A\mid (\bigvee X)\vee x=1\}\\
        &=\bigvee\{x\in A\mid x'\vee x=1, x'\in X\}\\
        &=\bigvee \alpha_A[X],
        \end{split}
        \]
        donde $\{x\in A\mid x'\vee x=1, x'\in X\}$ dirigido.
        \end{enumerate}
        \item Por último, notemos que para un marco eficiente, el morfismo $\alpha_A$ es suprayectivo ya que si $F=\nabla(j)\in A^\wedge$, entonces
        \[
        v_F=u_d
        \]
        donde $d=v_F(0)\in A$, es decir, $\alpha_A(d)=\nabla(u_d)=\nabla(v_F)=\nabla(j)$
    \end{enumerate}
\end{proof}

\begin{cor}
    Consideremos $c\in \mathfrak{c}A$ tal que $u_c$ es ajustado, entonces $\alpha_A$ es inyectiva. 
\end{cor}

\begin{proof}
    Consideremos $\nabla(u_c)=\nabla(u_d)$, con $d\in \mathfrak{c}A$. De aquí que $\nabla(u_c)\subseteq \nabla(u_d)$ y $\nabla(u_c)\supseteq\nabla(u_d)$. Aplicando dos veces el Lema \ref{cn ajustado} tenemos que 
    \[
    u_c\leq u_d\quad\mbox{ y }\quad u_c\geq u_d.
    \]
    Por lo tanto, evaluando ambos núcleos en $0$ obtenemos $c=d$.
\end{proof}

Notemos que el morfismo $\alpha_A$ no siempre es suprayectivo, pero existen maneras de asegurar que esto ocurra.

\begin{lem}\label{Lema2.7}
    Para un marco eficiente y $c\in \mathfrak{c}A$ las siguientes son equivalentes
    \begin{enumerate}
        \item $\alpha_A$ es un isomorfismo.
        \item Si $u_c\in \mathfrak{C}A$, $u_c$ es ajustado.
        \item Para todo $c\in \mathfrak{c}A$ y $d, d'\in A_{u_c}$
        \begin{equation}\label{saju}
            d\leq d'\Leftrightarrow \nabla(u_d)\subseteq \nabla(u_{d'}).
        \end{equation}
    \end{enumerate}
\end{lem}

\begin{proof}
\begin{description}
        \item[$1)\Rightarrow 2)$] Si $c\in \mathfrak{c}A$, entonces $u_c\in NA$ y tomemos $F=\nabla(u_c)\in A^\wedge$. Sea $v_F$ el respectivo núcleo ajustado del bloque, es decir, $v_F\leq u_c$. Además, por la eficiencia $v_F=u_d$ y  $\nabla(u_d)=\nabla(u_c)$. Al ser $\alpha_A$ inyectiva, $d=c$. Por lo tanto $u_c$ es ajustado. 

        \item[$2)\Rightarrow 3)$] Supongamos que $\nabla(u_c)\in A^\wedge$ y que $u_c$ es ajustado para $c\in \mathfrak{c}A$. Notemos que la implicación $\Rightarrow)$ de (\ref{saju}) siempre es cierta, pues si $d\leq d'$ y $d\vee e=1$, entonces $d'\vee e=1$. Así solo basta probar la otra implicación.\\

        Queremos demostrar que $\nabla(u_d)\subseteq \nabla(u_{d'})$ implica que $d\leq d'$. Por el Lema \ref{cn ajustado}
        \[
        \nabla(u_d)\subseteq \nabla(u_{d'})\Leftrightarrow u_d\leq u_{d'}\Leftrightarrow d\leq d'.
        \]

        \item[$3)\Rightarrow 1)$] Si $A$ es un marco es eficiente, por el Lema \ref{Lema2.8y2.9}, se cumple que $\alpha_A$ es suprayectiva. Para la inyectividad, consideremos $\alpha_A(d)=\alpha_A(d')$, es decir, $\nabla(u_d)=\nabla(u_{d'})$. Luego, aplicando dos veces (\ref{saju}), se cumple que $d=d'$. Por lo tanto, $\alpha_A$ es inyectiva.
    \end{description}
\end{proof}

Los siguientes resultados son consecuencias del Lema \ref{Lema2.7}.

\begin{cor}
        Para un marco $KC$ las siguientes son equivalentes
    \begin{enumerate}
        \item $\alpha_A$ es un isomorfismo.
        \item Si $j\in \mathfrak{K}A$, $j$ es ajustado.
        \item Para todo $c\in \mathfrak{c}A$ y $d, d'\in A_{u_c}$
        \[
            d\leq d'\Leftrightarrow \nabla(u_d)\subseteq \nabla(u_{d'}).
        \]
    \end{enumerate}
\end{cor}

\begin{cor}
    Para todo $c\in \mathfrak{c}A$, $u_c$ es ajustado si y solo si $\alpha_A$ es isomorfismo y $A$ es eficiente.
\end{cor}

\begin{obs}
Si $A$ es un marco eficiente, $j\in \mathfrak{K}A$ y $F=\nabla(j)\in A^\wedge$, entonces $j\in [v_F, w_F]$. Por la eficiencia  
\[
u_d=v_F\leq j\leq w_F\Rightarrow \nabla(u_d)=\nabla(j),
\]
es decir, si consideremos $j\in \mathfrak{K}A$ le asignamos un elemento $d\in \mathfrak{c}A$ a través de
\[
\varphi\colon\mathfrak{K}A\to \mathfrak{C}A
\]
donde $\varphi(j)=u_d$.

\end{obs}

\textbf{PREGUNTA:} ¿Podemos caracterizar a los marcos eficientes por medio de las funciones $\alpha_A$ y $\varphi$?\\

Para $\mathfrak{Q}S=\{j\in \mathfrak{K}A\mid j \mbox{ es ajustado}\}$ tenemos 
\[\begin{tikzcd}
	{\mathfrak{K}A} && {\mathfrak{C}A} & {\mbox{ y }} & {\mathfrak{K}A} && {\mathfrak{Q}A}
	\arrow["\varphi", shift left=3, from=1-1, to=1-3]
	\arrow["i", shift left=3, from=1-3, to=1-1]
	\arrow["{\varphi'}", shift left=3, from=1-5, to=1-7]
	\arrow["\iota", shift left=3, from=1-7, to=1-5]
\end{tikzcd}\]
donde $\iota\colon \mathfrak{C}A\to \mathfrak{K}A$ es una inclusión.\\

Notemos que $\varphi$ y $\varphi'$ son morfismos que desinflan pues $\varphi(k),\varphi'(k)\leq k$. Además, son monótonos, ya que para $j, k\in \mathfrak{K}A$ con $j\leq k$, se cumple que $u_d=\varphi(j)\leq j\leq k$. De aquí que $\nabla(u_d)\subseteq\nabla(k)=\nabla(u_{d'})$, donde $u_{d'}=\varphi(k)$. Como $u_d$ es un núcleo ajustado
\[
\nabla(u_d)\subseteq \nabla(u_{d'})\Rightarrow u_d\leq u_{d'}.
\]
El mismo argumento se aplica para $\varphi'$.\\

Verifiquemos que $\varphi\dashv \iota$, es decir, para $j\in \mathfrak{K}A$ y $u_c\in \mathfrak{C}S$ se cumple que 
\[
\varphi(j)\leq u_c\Leftrightarrow j\leq \iota(u_c).
\]
La implicación $\Leftarrow)$ es trivial. Para la otra implicación, Supongamos que $\varphi(j)\leq u_c$ y $j\nleq \iota(u_c)=u_c$. Como $u_d$ es ajustado, es el menor elemento de su bloque y como $j\nleq u_c$. Tenemos dos opciones
\[
i)\;  u_d\leq u_c\leq j\leq w_F \quad \mbox{o}\quad ii)\;u_d\leq u_c\; \mbox{ y }\;u_c\notin [u_d, w_F].
\]
Si ocurre $i)$, entonces 
Además, si el marco $A$ es eficiente $\mathfrak{C}A\simeq \mathfrak{Q}A$. Entonces tenemos el diagrama 

\[
\begin{tikzcd}
\mathfrak{C}A \arrow[r, "i"] \arrow[rr, "\alpha", bend left] & \mathfrak{K}A \arrow[r, "\hat{\alpha}"] \arrow[l, "\varphi", bend left] & A^\wedge
\end{tikzcd}
\]
donde $\hat{\alpha}(j)=\nabla(j)$


\subsection{Algunos núcleos de $N\mathfrak{c}A$} 

Consideremos $A\in \Frm$ y $a, b\in A$. Sabemos que $a\prec b$ si y solo si existe $x\in A$ tal que $a\wedge x=0$ y $c\vee c=1$. En otras palabras, para $x=\neg a$, $x\vee b=1$.

\begin{dfn}\label{Bb en CA}
    Para $c\in \mathfrak{c}A$ y $a,b\in A_{u_c}$ decimos que $a\prec_c b$ para tener la relación \emph{bastante por debajo} en el marco $A_{u_c}$.
\end{dfn}

El elemento $(a\succ c)$ es la negación de $a$ en $A_{u_c}$. De aquí que
\[
a\prec_cb\Leftrightarrow b\vee (a\succ c)=1.
\]

Si $c\leq d$, entonces 
\[
A_{u_c}=\uparrow c\supseteq \uparrow d=A_{u_d}
\]
produce un morfismo de marcos $i_{dc}\colon A_{u_d}\to A_{u_c}$ dado por $i_{dc}(x)=d\vee x$. De esta manera para $c\leq d\leq e$ tenemos
\[\begin{tikzcd}
	{A_{u_e}} && {A_{u_d}} \\
	& {A_{u_c}}
	\arrow["{i_{ed}}"', from=1-3, to=1-1]
	\arrow["{i_{ec}}", from=2-2, to=1-1]
	\arrow["{i_{dc}}"', from=2-2, to=1-3]
\end{tikzcd}\]
y $i_{cc}\colon A_{u_c}\to A_{u_c}$ es la identidad.\\

La construcción anterior produce un funtor $F\colon \mathfrak{c}A\to \Frm$ donde 
\[
c\mapsto A_{u_c}\quad \mbox{ y }\quad c\leq d\mapsto i_{dc}
\]
son la asignación en objetos y flechas, respectivamente.\\

Consideremos un cono sobre $F$ dado por $(\hat{A}, \{\pi_c\colon \hat{A}\to A_{u_c}\}_{c\in \mathfrak{c}A})$. De manera similar que en la categoría de conjuntos, 
\[
\prod_{c\in \mathfrak{c}A}F(c)=\prod_{c\in \mathfrak{c}A}A_{u_c}=\{j\colon \mathfrak{c}A\to \bigcup A_{u_c}\mid j(c)\in A_{u_c}\}
\]
para todo $c\in \mathfrak{c}A$ y $j$ una función. De aquí que $j(c)\in A_{u_c}$ si y solo si $c\leq j(c)$, es decir, 
\[
j\in \prod_{c\in \mathfrak{c}A}A_{u_c} \Leftrightarrow c\leq j(c).
\]

Además, si $c\leq j(c)$ entonces $\pi_c(j)=j(c)=c\vee j(c)$, es decir, para $c\leq d$, $j(d)=i_{dc}(j(c))=d\vee j(c)$ y $j\colon \mathfrak{c}A\to \mathfrak{c}A$. Por lo tanto
\begin{equation}\label{compatibilidad}
\begin{split}
\hat{A}&=\{j\in \prod_{c\in \mathfrak{c}A}A_{u_c}\mid j(d)=i_{dc}(j(c))\}\\
&=\{j\colon \mathfrak{c}A\to \mathfrak{c}A\mid j(c)\in A_{u_c}\}
\end{split}
\end{equation}
para todo $c\leq d$. 

\begin{prop}\label{Flimite}
    Si $c\leq d\in \mathfrak{c}A$, entonces $\hat{A}$ es un cono límite sobre $F$.
\end{prop}

\begin{proof}
\begin{enumerate}
    \item Por construcción, cada proyección $\pi_c$ está bien definida. Además, si $c\leq d$
    \[
    i_{dc}\circ \pi_c(j)=i_{dc}(j(c))= d\vee j(c)=j(d)= \pi_d(j),
    \]
es decir, los triángulos del cono conmutan. Por lo tanto, $\hat{A}$ es un cono sobre $F$.

\item Sea $(Y,\{f_c\colon Y\to A_{u_c}\}_{c\in \mathfrak{c}A})$ cualquier otro cono sobre $F$, esto es,
para todo $c\le d$ se cumple $i_{dc}\circ f_c=f_d$. 
Consideremos el morfismo $u:Y\to \hat{A}$ definido por
\[
(u(y))(c)=f_c(y)\qquad (y\in Y,\ c\in \mathfrak{c}A).
\]
Primero vemos que $u(y)\in\lim F$ para cada $y\in Y$. Si $c\le d$, entonces
\[
(u(y))(d)=f_d(y)=(i_{dc}\circ f_c)(y)=i_{dc}(f_c(y))
=d\vee(u(y))(c),
\]
y por (\ref{compatibilidad}), $u(y)\in \hat{A}.$

Además, por definición de $u$ y de las proyecciones, para cada $c$ se tiene
\[
\pi_c\circ u \;=\; f_c,
\]
de modo que $u$ es un morfismo de conos $(Y\to \hat{A})$. Veamos que $u$ es único.\\

Si $v\colon Y\to \hat{A}$ es otro morfismo de conos con 
$\pi_c\circ v = f_c$ para todo $c$, entonces para todo $y\in Y$ y $c\in\mathfrak{c}A$,
\[
(v(y))(c)=\pi_c(v(y))=f_c(y)=(u(y))(c).
\]
De aquí que $v(y)=u(y)$ para todo $y$, luego $v=u$.
\end{enumerate}
Por lo tanto $\hat{A}$ satisface la propiedad universal del cono límite.
\end{proof}

\begin{lem}\label{Lema3.1}
    Para cualquier $A\in \Frm$ y cada función $j\colon \mathfrak{c}A\to \mathfrak{c}A$ las siguientes son equivalentes:
    \begin{enumerate}
        \item $j\in \hat{A}$.
        \item $j(d)=j(c)\vee d$ para todo $d\geq c\in \mathfrak{c}A$.
        \item $j(c\vee a)=j(c)\vee a$ para todo $c\in \mathfrak{c}A$ y $a\in A$.
        \item $j(c\vee e)=j(c)\vee e$ para todo $c, e\in \mathfrak{c}A$.
    \end{enumerate}
\end{lem}

\begin{proof}
    \begin{description}
        \item[$1)\Leftrightarrow 2)$] Se cumple por la construcción de $\hat{A}$.
        \item[$2)\Rightarrow 3)$] Tomando $c\leq d=c\vee a$, $d\in \mathfrak{c}A$, pues $\mathfrak{c}A$ es una sección superior. De esta manera, por $2)$
        \[
        j(d)=j(c)\vee d=j(c)\vee c\vee a=j(c)\vee a,
        \]
        pues $c\leq j(c)$.
        \item[$3)\Rightarrow 4)$] Si $e\in \mathfrak{c}A$, $e\in A$ y aplicando $3)$ obtenemos lo que queremos.
        \item[$4)\Rightarrow 2)$] Si $c\leq d$, se cumple que $c\vee d=d$, tomando $e=d$ y aplicando $4)$ se tiene $j(d)=j(c)\vee d$. 
    \end{description}
\end{proof}

\begin{lem}\label{Lema3.3}
    Para $A\in \Frm$ si $j\in \hat{A}$, entonces $j$ es un núcleo en $\mathfrak{c}A$.
\end{lem}

\begin{proof}
    Verifiquemos que $j\in \hat{A}$ cumple las condiciones de núcleo.
    \begin{enumerate}
        \item Por construcción, $c\leq j(c)$, es decir, $j$ infla.
        \item Notemos que para $d=j(c)$, entonces $j(d)=j(j(c))=j(c)\vee j(c)=j(c)$. Por lo tanto $j$ es idempotente.
        \item Consideremos $c=d\wedge d'$, entonces $c\leq d, d'$. Luego
        \[
        j(d)=j(d\wedge d')\vee d\quad \mbox{ y }\quad j(d')=j(d\wedge d')\vee d'.
        \]
        De aquí que 
        \[
        \begin{split}
            j(d)\wedge j(d')& =(j(d\wedge d')\vee d)\wedge (j(d\wedge d')\vee d')\\
            & =j(d\wedge d')\vee (d\wedge d')\\
            &=j(d\wedge d')
        \end{split}
        \]
        donde la segunda igual se da por la distributividad y la última del hecho de que $j$ infla. Así $j$ respeta ínfimos.
    \end{enumerate}
    Por lo tanto, $j$ es un núcleo en $\mathfrak{c}A$.
\end{proof}

Sabemos que si $j\in NA$, entonces $j^{-1}(1)$ es un filtro.

\begin{lem}\label{Lema3.3}
    Para cualquier núcleo $j$ en una semiretícula de Heyting $A$ y cualquier $a\in A$, la desigualdad $v_a\leq j$ se cumple si y solo si $j(a)=1$, donde $v_a(x)=(a\succ x)$.
\end{lem}

\begin{proof}
Es consecuencia de las propiedades de los $v$-núcleos.
\end{proof}

Enunciamos un lema auxiliar para la prueba del siguiente teorema.

\begin{lem}\label{lem:nucleo-dirigidos-a-arbitrarios}
Sea $A$ un marco y $j\colon A\to A$ un núcleo. 
Si $j$ preserva supremos finitos y supremos dirigidos, entonces $j$ preserva todos los supremos arbitrarios, es decir,
\[
j(\bigvee X)=\bigvee j[X]
\]
para todo $X\subseteq A$.
\end{lem}

\begin{proof}
Sea $X\subseteq A$ y consideremos el conjunto
\[
\mathcal{F}(X)=\{\bigvee F\mid F\subseteq X\text{ es finito}\}.
\]
Notemos que $\mathcal{F}(X)$ es dirigido pues $u=\bigvee F$ y $v=\bigvee G$ con $F,G$ finitos, 
entonces $w=\bigvee(F\cup G)\in\mathcal{F}(X)$ y $u\le w$, $v\le w$.

Además, 
\[
\bigvee X\;=\;\bigvee \mathcal{F}(X).
\]

Por hipótesis $j$ preserva supremos finitos y dirigidos, entonces
\[
j(\bigvee X)=j(\bigvee \mathcal{F}(X))=\bigvee_{u\in \mathcal{F}(X)} j(u)=\bigvee_{F\subseteq_{\mathrm{fin}} X} j(\bigvee F)=\bigvee_{F\subseteq_{\mathrm{fin}} X} \bigvee j[F].
\]
Finalmente, como $\{\bigvee j[F]\mid F\subseteq_{\mathrm{fin}} X\}$ es dirigido y su supremo coincide con $\bigvee j[X]$, tenemos que
\[
j\!\Big(\bigvee X\Big)\;=\;\bigvee j[X].
\]
\end{proof}

\begin{thm}\label{Teorema3.4}
    Para $A$ un marco Hausdorff y cualquier núcleo $j\colon \mathfrak{c}A\to \mathfrak{c}A$ las siguientes son equivalentes
    \begin{enumerate}
        \item $j\in \hat{A}$.
        \item $j$ preserva supremos no vacíos.
        \item $j$ es Scott-continuo.
        \item $\nabla(j)\in \mathfrak{c}A^\wedge$.
    \end{enumerate}
\end{thm}

\begin{proof}
    \begin{description}
        \item[$1)\Rightarrow 2)$] Consideremos $c, d\in \mathfrak{c}A$, entonces
        \[
        j(c\vee d)=j(c)\vee c\vee d=j(c)\vee d\leq j(c)\vee j(d).
        \]
        Por monotonía se cumple que $j(c)\vee j(d)\leq j(c\vee d)$. Por lo tanto $j(c\vee d)=j(c)\vee j(d).$

        Para los supremos dirigidos, consideremos $\mathcal{D}\subseteq \mathfrak{c}A$ dirigido y tomemos $c\in \mathcal{D}$. Como $\bigvee \mathcal{D}\in \mathfrak{c}A$ se cumple que 
        \[
        j(\bigvee \mathcal{D})=j(c\vee \bigvee \mathcal{D})=j(c)\vee\bigvee \mathcal{D}=\bigvee_{d\in \mathcal{D}}j(c)\vee d=\bigvee_{d\in \mathcal{D}}j(c\vee d)=\bigvee_{e\in \mathcal{D}}j(e)
        \]
        donde la última igual se cumple por ser $\mathcal{D}$ dirigido. 

        De aquí que, por el Lema \ref{lem:nucleo-dirigidos-a-arbitrarios} se cumple que $j$ preserva cualquier supremo.

        \item[$2)\Rightarrow 3)$] Si $j$ preserva supremos dirigidos, $j$ es Scott continuo.

        \item[$3)\Rightarrow 4)$] Por la prueba del Lema \ref{Lema2.3} $5)$, como $\mathfrak{c}A$ es compacto se tiene que para $X\subseteq \mathfrak{c}A$, con $X$ dirigido, existe $x\in X$ tal que $x=1$ y $x\in A_{u_c}$ para $c\in \mathfrak{c}A$ y $1=\bigvee X$. como $j$ es Scott continuo
        \[
        j(\bigvee X)=\bigvee j[X]=1,
        \]
        es decir, $\bigvee X\in \nabla(j)$. de aquí que existe $j(x)\in j[X]$ tal que $j(x)=1$, entonces $x\in \nabla(j)$. Por lo tanto, $X\cap\nabla(j)\neq\emptyset$ y $\nabla(j)\in \mathfrak{c}A^\wedge$.

        \item[$4)\Rightarrow 1)$] Sabemos que $j\in \lim F$ si $j(d)=j(c)\vee d$ para $c\leq d$. Como $j$ es monótona
        \[
        j(c)\leq j(d)\Rightarrow d\vee j(c)\leq d\vee j(d)=j(d)
        \]
        siempre se cumple. De esta manera, basta con verificar que $j(d)\leq d\vee j(c)$. Consideremos $e\in A_{u_c}$ tal que $e\vee j(d)=1$ y como $e\vee j(d)\leq j(e\vee d)$, entonces $j(e\vee d)=1$, es decir, $e\vee d\in \nabla(j)$.
    \end{description}
\end{proof}

\section{Preguntas abiertas}

\begin{enumerate}
    \item Bajo las siguientes asignaciones 
    \[
    \nabla(\_ )\colon Fil(A)\to A^\wedge\quad \iota\colon A^\wedge \to Fil(A)
    \]
    ¿Pasa qué $\nabla(\_ ) \vdash \iota$? esto es cierto solo falta hacer la cuenta
    
    \item ¿Si un marco es subajustado produce cocientes compactos?
          %por que queriamos esto?%
    \item ¿Cuándo el intervalo $[v_\nabla (0), w_\nabla(0)]$ colapsa?

    \item Para $V_{v_\nabla}\wedge U_{w_\nabla}\in N^2A$ y $V_{v_\nabla(0)}\wedge U_{w_\nabla(0)}\in NA$, ¿qué relación existe entre estos núcleos?

    \item Consideremos el conjunto $\{[v_F, w_F]\mid F\in A^\wedge\}\subseteq I(NA)$, ¿qué propiedades satisface este conjunto?

    \item ¿Se puede estratificar la noción de ajustado?
  %habia unas cosas con ciertas categorias de marcos no?%

  \item ¿Es todo marco Hausdorff parche-trivial, es decir, es todo marco Hausdorff arreglado?
  
  \item ¿Cómo podemos medir el fallo del arreglo de un marco ?

  \item ¿$PPPA=PPA$?
  
\item ¿Qué propiedades de separación satisface el marco de parches de cualquier marco? Recordar que $N(A)$ siempre es regular puesto que es cero dimensional.

\item Si $k\in NA$ es continuo, ¿qué sabemos sobre $\nabla(k)$?

\item ¿Qué significa la noción de apilado para marcos?

\item El cuadrado que relaciona un marco $A$ con la topología de su espacio de puntos conmuta laxamente en $\mathbf{Pos}$, es decir, $U_Av_F\leq \nabla^\infty U_A$. ¿Cuándo se da la igualdad? 

\item ¿Cuál es el comportamiento del $Q$-cuadrado y los intervalos $[v_\_, w_\_]$ cuando se supone alguna propiedad tipo Hausdorff

\end{enumerate}

\end{document}
